\section{Analysis and Discussion}

The results support the paper's central claim: privacy containment in AI-mediated platforms requires mediating intermediate state transitions. Sensitive content may first appear as a workspace write, summary, or inter-agent message; once that artifact becomes readable, later agents can transform it into routine operational text. Final-output filters observe only the end of this chain. FlowFence acts earlier, at the transition into shared state or external channels.

The topology result clarifies why this is a platform externality. A local write has different downstream costs depending on where it lands. Private memory, a one-hop chain message, a planner-centered star, and a shared blackboard create different propagation surfaces. Static ACLs remain useful for hard forbidden channels, but they do not capture provenance, privilege transitions, fanout, or abstraction rights. Prompt filters are complementary but depend on surface form. FlowFence instead evaluates risk from policy, provenance, topology, privilege, and runtime content signals.

The combination of task success and zero unauthorized raw/external leakage in the evaluated FlowFence runs rules out a trivial block-all explanation. The vendor-facing task still completes because safe views preserve authorized abstractions, while quarantine prevents raw protected values from entering high-fanout shared state. This distinction is important for platform governance: useful coordination often requires disclosing that a constraint exists, not disclosing the raw sensitive value.

The label-free validation reduces the main internal-validity concern. Removing evaluator-only attack annotations from guard input and testing reserved paraphrase variants shows that the reported containment is not simply an artifact of attack labels. The result does not imply coverage of every semantic paraphrase or adaptive adversary; it supports the narrower claim that runtime-observable signals can contain the evaluated propagation patterns.
